\documentclass{article}
\usepackage[T1]{fontenc}
\usepackage{lmodern}
\usepackage{graphicx}
\usepackage{amsmath,amssymb}
\usepackage[a4paper,margin=2cm]{geometry}
\usepackage[absolute,overlay]{textpos}
\setlength{\TPHorizModule}{1mm}
\setlength{\TPVertModule}{1mm}
\usepackage[indonesian]{babel}
\usepackage{hyperref}
\setlength{\emergencystretch}{2em}
% unit: Halaman 4

\begin{document}

% === Halaman 4 (python/native) ===

Lima finalis lomba AES:

1. Rijndael (dari Vincent Rijmen dan Joan Daemen – Belgia, 86 suara)

2. Serpent (dari Ross Anderson, Eli Biham, dan Lars Knudsen – 

Inggris, Israel, dan Norwegia, 59 suara).

3. Twofish (dari tim yang diketuai oleh Bruce Schneier – USA, 31 

suara)

4. RC6 (dari Laboratorium RSA – USA, 23 suara)

5. MARS (dari IBM, 13 suara)

\end{document}
